% Рисунок с доски: поле направлений a(x) и его интегральные кривые.
% Решение уравнения a.grad u = 0 постоянно вдоль каждой такой кривой.
\begin{figure}[htbp]
\centering
\begin{tikzpicture}[>={Latex[length=4pt]}, scale=1.0]
  \draw[->] (-0.3,0) -- (6.2,0) node[anchor=north east] {$x_1$};
  \draw[->] (0,-0.3) -- (0,3.6) node[anchor=south east] {$x_2$};

  % семейство интегральных кривых
  \foreach \k/\o in {0.9/0.0, 0.9/0.55, 0.9/1.1, 0.9/1.65} {
    \draw[thick, blue!60!black]
      plot[domain=0.25:5.6, samples=60]
      (\x, {\o + \k*sqrt(\x)});
  }

  % касательные векторы
  \draw[->, red!65!black, thick] (1.6,1.138) -- (2.5,1.44);
  \node[red!65!black, anchor=north west, font=\footnotesize] at (2.3,1.42)
    {$\bar a$};
  \draw[->, red!65!black, thick] (3.2,2.26) -- (4.1,2.47);
  \node[red!65!black, anchor=south east, font=\footnotesize] at (4.1,2.47)
    {$\bar a$};
  \draw[->, red!65!black, thick] (0.9,0.854) -- (1.8,1.3);
  \node[red!65!black, anchor=north west, font=\footnotesize] at (1.55,1.26)
    {$\bar a$};
\end{tikzpicture}
\caption{Поле направлений $\bar a(x)$ и его интегральные кривые
(характеристики). Уравнение $\bar a(x)\cdot\nabla u = 0$ означает, что
производная $u$ вдоль каждой такой кривой равна нулю: решение постоянно на
характеристиках.}
\label{fig:char-field}
\end{figure}
